\documentclass[a4paper]{article}
\usepackage[pdftex]{graphicx}
\usepackage[utf8]{inputenc}
\usepackage{enumerate}
\usepackage{icomma}
\usepackage{amssymb}
\usepackage{tikz}
\usepackage{href-ul}
\hypersetup{
	colorlinks=true,
	linkcolor=blue,
	urlcolor=blue}
\usepackage{geometry}
\geometry{a4paper, top=15mm, left=15mm, right=15mm, bottom=15mm,
	headsep=10mm, footskip=12mm}

\begin{document}
\begin{tikzpicture}(10,3)
	\draw[ultra thick](2,0) --(10,0) -- (10,1.3) --(2,1.3) -- (2,0);
	\draw[fill=black](2,0)-- (0,.6) -- (2,.6) -- (2,0);
	\node at (6,.6) {\large Lösungsblatt von \href{https://www.okuyakl.de}{www.okuyakl.de}};
\end{tikzpicture}
\vspace{0.5 cm}

\noindent{\bf Aufgabe 1. a)}\\
Wir können ausklammern und erhalten als doppelte Nullstelle:
$$x^2 ({3\over4 }x^2-4x +6)=0 \quad \Rightarrow \quad x_{1/2}=0$$
Untersuchung des quadratischen Terms in der Klammer liefert für die Diskriminante $D$:
$$D=b^2-4ac=4^2-4\cdot {3\over 4}\cdot 6 = -2 < 0$$
Es existieren keine weiteren Nullstellen.

\noindent{\bf Aufgabe 1. b)}\\
Wir leiten ab und setzen gleich Null:
$$
\renewcommand{\arraystretch}{2}
\begin{array}{rcll}
f'(x)=0&=&3x^3-12x^2+12x \\
      0&=&3x(x^2-4x+4) & x_1=0 \\
      0&=&(x-2)^2 & x_{2/3}=2
\end{array}
$$ 
Die Koordinaten der gefundenen Stellen erhalten wir durch Einsetzen in $f(x)$:
$$f(0)=0 \quad \Rightarrow \quad P_1(0|0)$$
$$f(2)={3\over 4}\cdot 16 - 4\cdot 8 +6\cdot 4 = 4 \quad \Rightarrow \quad P_2(2|4)$$

\noindent{\bf Aufgabe 2.}\\
Allgemeiner Funktionsterm 3. Grades:
$$f(x)=ax^3+bx^2+cx+d$$
Eine ganzrationale Funktion ist punktsymmetrisch, wenn sie nur ungerade Potenzen von $X$ enthält:
$$\Rightarrow \quad f(x)=ax^3+cx$$
Sie enthält den Punkt $H(-1|4)$:
$$f(-1)=-a-c=4 \quad I $$
Am Hochpunkt ist die Ableitung null:
$$f'(x)=3ax^2 + c \quad \Rightarrow \quad 3a +c = 0 \quad II $$
Wir setzen I in II ein und erhalten:
$$ 3a -4 -a =0 \quad \Rightarrow \quad a = 2  \quad \Rightarrow \quad c = -6 $$
Der Funktionsterm lautet:
$$f(x)= 2x^3 -6x$$

\noindent{\bf Aufgabe 3. a)}\\
$$f'(x)=3x^5-2,5 + 4 \cdot x^{-3}$$

\noindent{\bf Aufgabe 3. b)}\\
$$f'(x)=(2x^{2,5})' = 5x^{1,5}$$

\noindent{\bf Aufgabe 4. a)}\\
Wir leiten ab:
$$f'(x)= 0,75 x^2 + 6x +9 $$
Die Tangentensteigung im Punkt $B(-4|-4)$ erhalten wir durch Einsetzen:
$$ m = f'(-4)= 12 - 24 +9 = -3$$
Der gesuchte Winkel erfüllt die Beziehung:
$$ m = \tan\alpha \quad \Rightarrow \quad \alpha =\tan^{-1}(-3) = -71^\circ = 109^\circ$$

\noindent{\bf Aufgabe 4. b)}\\
Die Funktion hat zwei Tangenten dieser Steigung, wenn es zwei Lösungen der Gleichung gibt:
$$
\renewcommand{\arraystretch}{2}
\begin{array}{rcll}
 0,75 x^2 + 6x +9 &=& -3 \\
 0,75 x^2 +6x +12 &=& 0 \\
 x_{1/2} &=& {-6 \pm \sqrt{36-4\cdot 0,75 \cdot 12}\over 1,5}\\
 x_1 = x_2 &=& -{6 \over 1,5}= - 4 
\end{array}
$$
Es gibt nur eine Lösung; also gibt es keine weitere parallele Tangente. (Für Interessierte: die Tangente geht durch den Wendepunkt)

\noindent{\bf Aufgabe 5.}\\
Zunächst leiten wir ab und setzen die beiden Nullstellen ein:
$$f'(x)=-2x + 5 \quad \Rightarrow \quad f'(0)= 5 \qquad f'(5)=-5$$ 

Die zugehörigen Normalensteigungen sind die negativen Kehrwerte davon:
$$m_1 = -{1 \over 5} \qquad m_2 = {1\over 5}$$

Damit sind die beiden Funktionsterme der Normalen:
$$ 0 = t_1 \quad \Rightarrow \quad n_1: y= -{1 \over 5}x $$
$$ 0 = {1\over 5} \cdot 5 + t_2 \quad \Rightarrow \quad t_2 = -1 \quad \Rightarrow \quad n_2: y={1\over 5}x-1 $$

Wir setzen beide Geraden gleich und erhalten für den Schnittpunkt:
$$ -{1 \over 5}x = {1\over 5}x-1  \quad \Rightarrow \quad x_S = 2,5  \quad \Rightarrow \quad y_S = -0,5 \quad \Rightarrow \quad P(2,5|-0,5)$$

Das gesuchte Dreieck ist gleichschenklig mit der Basis auf der x-Achse, die Basislänge ist 5. Die Höhe des Dreiecks ist die y-Koordinate von P, also -0,5. Die Dreiecksfläche ist der Betrag davon:
$$ A = {1\over 2} \cdot 5 \cdot 0,5 = 1,25~FE$$

\noindent{\bf Aufgabe 6. a)}\\
Ein möglicher Funktionsterm lautet:
$$f(x)=-x^2+9 $$



\begin{minipage}{0.5\textwidth}
\noindent{\bf Aufgabe 6. b)}\\
\includegraphics[width= 7 cm]{pavier060}
\end{minipage}
\begin{minipage}{0.45\textwidth}
\noindent{\bf Aufgabe 6. c)}\\
Das Rechteck hat die Breite $2u$ und die Höhe:
$$f(u)=-u^2+9$$
Die Fläche ist das Produkt aus beiden, also:
$$A(u)=2u\cdot(-u^2+9) = -2u^3 + 18u $$
Um das Maximum zu bekommen, leiten wir ab und setzen gleich null:
$$A'(u)=-6u^2+18 =0 \Rightarrow u_{1/2}=\pm \sqrt{3}$$
Für $u=\sqrt{3}$ wird die Fläche maximal. 
\end{minipage}

\begin{center}
	\includegraphics[width=7 cm]{../../viecher/endcomic.pdf}
	
	Hier geht es zurück zum \href{https://www.okuyakl.de/math/m11garaL060/aa060.pdf}{Aufgabenblatt}
\end{center}


\end{document}

